% Pursuing Stacks (cotes 134-2 to 134-8): archive cote 134-7, site page 23 (PDF page 24), typescript page 465.
% typescript. Draft compared: ps6.tex:1066-1125.
% AI-transcribed by Claude Opus 5.5 (claude-opus-5-5, adaptive thinking, effort high), Message Batches API, batch ps3 (msgbatch_01XdFmuVVBAcgc7EFxdtCqsj), 2026-11-12; not yet checked by eye.
\scanpage{134-7/23}
\mspage{465}

If it is possible indeed to construct Postnikov dévissage of a
\xout{\unsure{Kan}} schematic Kan complex over any ground ring k, it is clear that this
is compatible with restriction of ground ring, hence it would seem that
formation of the homotopy \xout{\unsure{module}} invariants $\hand{\pi}_i$ is compatible with \xout{\illegible{1}}
restriction of rings\deleted{\unsure{,}} (whereas, as we noticed yesterday, the same does
definitely \uline{not} hold for the homology invariants H$_i$). Taking restriction
to the absolute ground ring $\underline{Z}$, this shows that the canonical functor
(***) from schematic to discrete homotopy types is compatible with
taking homotopy groups (but not with homology) - thus the relation
between $X_*$ and the complex of sections $X_*(k)$ seems to be a rather close
one, via the homotopy groups\hand{,} which are the same (and thus, the homotopy
groups of $X_*(k)$ seem to turn out to be k-modules after all !). By the
way, speaking of ``restriction of ground ring'' for Kan complexes was a
little hasty, in view of the projectivity condition on the components,
which a priori \inserted{seems to} oblige\deleted{s} us to assume k to be a projective k$_o$-module
(for a given ring homomorphism
\[ k_o \hand{\longrightarrow} k \quad ) . \]
Still, the remark about the sections functor $X_* \hand{\longrightarrow} X_*(k)$ makes a
sense, without having to assume k to be a projective $\underline{Z}$-module ! Also,
we feel that, by analogy of what can be done in the linear set up,
when we define a total derived functor
\[ D.(Ab_k) \hand{\longrightarrow} D.(Ab_{k_o}) \]
without any assumption on the ring homomorphism k$_o$ \hand{$\to$} k , a notion of
ring restriction for \inserted[typed]{schematic} homotopy types should make sense without any
restriction, as was surmised yesterday. As in the linear case, we
should allow ourselves to work with schematic complexes which are \uline{not}
projective, but be prepared to take ``resolutions'' \inserted{(in some sense)} of such general comple-%
xes by the more restricted ones. \hand{(with projective components).}

There is no such difficulty in the case of the ring extension func-%
tor, which transforms projective bundles over k$_o$ into projective bundles.
The reflexions above suggest that, whereas ring extension is compatible
with taking total homology invariants LH. , via the corresponding
functor
\[ D.(Ab_{k_o}) \hand{\longrightarrow} D.(Ab_k) \quad , \]
it is compatible too with taking homotopy invariants $\hand{\pi}_i$ separately.

\margin{23}
